\subsection{有限素数支撑的 Boolean exact calculus 与一轴 prime-register 的有限支撑骨架}\label{subsec:conclusion-finite-primesupport-boolean-exact-uniaxial-register}
沿用定理
\ref{thm:killo-localization-circle-dimension-prime-ledger-splitting}、
\ref{thm:conclusion-finite-prime-solenoid-terminal-object}、
\ref{thm:conclusion-finite-prime-localization-pontryagin-rigidity}、
\ref{thm:xi-localized-phase-sampling-closed-form}、
\ref{thm:killo-godel-semidir-affine-2adic}
与猜想
\ref{conj:conclusion-uniaxial-prime-register-universality}
的记号。前述结果已经分别给出：有限素数局部化的 Pontryagin 对偶短正合列、有限 prime-support completion 的终对象性与刚性分类、局部化相位采样对有限素数支撑的唯一恢复，以及 canonical finite-time-depth prime-register 在固定二维 \(2\)-进 Tate 模块上的忠实仿射实现。把这些接口并读之后，可进一步把有限素数支撑提升为一套在 compact--profinite 侧可做精确商、精确交与精确粘接的 Boolean exact calculus。

\begin{theorem}[有限素数支撑的精确商]\label{thm:conclusion-finite-primesupport-exact-quotient}
\leanverified{paper\_conclusion\_finite\_primesupport\_exact\_quotient}
对任意有限素数集
\[
S\subseteq T\subset \mathbb P,
\qquad
G_S:=\ZZ[S^{-1}],
\qquad
\Sigma_S:=\widehat{G_S},
\]
存在自然短正合列
\[
0\longrightarrow \prod_{p\in T\setminus S}\ZZ_p
\longrightarrow \Sigma_T
\xrightarrow{\ q_{T\to S}\ }
\Sigma_S
\longrightarrow 0.
\]
特别地，删去有限素数支撑中的每一个新增素数 \(p\in T\setminus S\)，在对偶侧恰对应于剥去一个 \(p\)-进因子 \(\ZZ_p\)，而连通圆因子保持不变。
\end{theorem}
\begin{proof}
由定理 \ref{thm:killo-localization-circle-dimension-prime-ledger-splitting}，对任意有限 \(U\subset \mathbb P\) 都有
\[
0\longrightarrow \ZZ\longrightarrow G_U\longrightarrow \bigoplus_{p\in U}C_{p^\infty}\longrightarrow 0
\]
及其 Pontryagin 对偶
\[
0\longrightarrow \prod_{p\in U}\ZZ_p\longrightarrow \Sigma_U\longrightarrow \TT\longrightarrow 0.
\]
当 \(S\subseteq T\) 时，离散侧有自然嵌入
\[
G_S=\ZZ[S^{-1}]\hookrightarrow \ZZ[T^{-1}]=G_T,
\]
并且商群恰为
\[
G_T/G_S\cong \bigoplus_{p\in T\setminus S}C_{p^\infty},
\]
因为新增的分母自由度正好来自 \(T\setminus S\) 中各素数的 Prüfer 塔。对该短正合列取 Pontryagin 对偶，便得到
\[
0\longrightarrow \prod_{p\in T\setminus S}\ZZ_p
\longrightarrow \Sigma_T
\xrightarrow{\ q_{T\to S}\ }
\Sigma_S
\longrightarrow 0.
\]
最后，由定理 \ref{thm:conclusion-finite-prime-solenoid-terminal-object}，每个 \(\Sigma_U\) 的非平凡连通 Lie 可见商都同构于 \(\TT\)，故上述商映射只剥去全不连通 prime-ledger 核，而不改变唯一的连通圆因子。证毕。
\end{proof}

\begin{theorem}[有限素数支撑的 Boolean--exact 双笛卡尔演算]\label{thm:conclusion-finite-primesupport-bicartesian-boolean-calculus}
\leanverified{paper\_conclusion\_finite\_primesupport\_bicartesian\_boolean\_calculus}
对任意有限素数集 \(S,T\subset \mathbb P\)，存在自然短正合列
\[
0\longrightarrow \Sigma_{S\cup T}
\longrightarrow \Sigma_S\times \Sigma_T
\xrightarrow{\ \delta\ }
\Sigma_{S\cap T}
\longrightarrow 0,
\]
其中
\[
\delta(x,y):=q_{S\to S\cap T}(x)-q_{T\to S\cap T}(y).
\]
因此，在 compact abelian groups 范畴中，方块
\[
\begin{matrix}
\Sigma_{S\cup T} &\longrightarrow& \Sigma_S\\
\downarrow && \downarrow\\
\Sigma_T &\longrightarrow& \Sigma_{S\cap T}
\end{matrix}
\]
同时是 pullback 与 pushout。换言之，有限素数支撑的并、交两种 Boolean 运算，在相位完成侧严格实现为一条双笛卡尔 exact calculus。
\end{theorem}
\begin{proof}
在离散侧考虑自然加法映射
\[
\psi:G_S\oplus G_T\longrightarrow G_{S\cup T},
\qquad
\psi(x,y):=x+y.
\]
先证其满射。任取
\[
z\in G_{S\cup T}.
\]
可将其写成
\[
z=\frac{a}{cmn},
\]
其中 \(c\) 的素数支撑包含于 \(S\cap T\)，\(m\) 的素数支撑包含于 \(S\setminus T\)，\(n\) 的素数支撑包含于 \(T\setminus S\)。于是 \((m,n)=1\)，故存在整数 \(\alpha,\beta\) 使
\[
\alpha m+\beta n=1.
\]
从而
\[
z=\frac{a}{cmn}
=
\frac{a\beta}{cm}+\frac{a\alpha}{cn},
\]
其中
\[
\frac{a\beta}{cm}\in G_S,
\qquad
\frac{a\alpha}{cn}\in G_T.
\]
故 \(\psi\) 满射。

其核恰为反对角嵌入的 \(G_{S\cap T}\)：若 \(x+y=0\)，则
\[
x=-y\in G_S\cap G_T=G_{S\cap T},
\]
反之任取 \(z\in G_{S\cap T}\)，对偶对 \((z,-z)\) 显然落在核中。故有短正合列
\[
0\longrightarrow G_{S\cap T}
\longrightarrow G_S\oplus G_T
\longrightarrow G_{S\cup T}
\longrightarrow 0.
\]
对偶化后得到
\[
0\longrightarrow \Sigma_{S\cup T}
\longrightarrow \Sigma_S\times \Sigma_T
\xrightarrow{\ \delta\ }
\Sigma_{S\cap T}
\longrightarrow 0.
\]
而 Pontryagin 对偶把离散交换群范畴与紧交换群范畴反变等价；在离散侧，此短正合列正对应于并/交的标准 Mayer--Vietoris 型演算，因此在对偶侧对应方块同时给出 fiber product 与 amalgamated pushout。证毕。
\end{proof}

\begin{theorem}[实现层二维与结构层一圆维的严格分离]\label{thm:conclusion-prime-register-twoadic-implementation-onecircle-structure-separation}
设
\[
H_{\mathrm{can}}
\]
为任意 canonical finite-time-depth prime-register 子单子。则同时成立：
\begin{enumerate}
  \item \(H_{\mathrm{can}}\) 可忠实实现为固定二维 \(2\)-进 Tate 模块
  \[
  T_2(E_{\min})\cong \ZZ_2^2
  \]
  上的仿射变换单子；
  \item 任意由局部化相位采样恢复的有限 prime-support 对象，其结构层只对应某个
  \[
  G_S=\ZZ[S^{-1}]
  \]
  及其对偶短正合列
  \[
  0\longrightarrow \prod_{p\in S}\ZZ_p\longrightarrow \Sigma_S\longrightarrow \TT\longrightarrow 0;
  \]
  \item 这些结构层对象恒满足
  \[
  \cdim_{\mathrm{fr}}(G_S)=1.
  \]
\end{enumerate}
因此，prime-register 机制在实现层上最小需要二维 \(2\)-进 ambient space，而在结构层上始终只携带唯一的连通圆因子；第二个 \(2\)-进方向只是 faithful realization 的工作坐标，而不是额外的结构圆维。
\end{theorem}
\begin{proof}
第一条正是定理 \ref{thm:killo-godel-semidir-affine-2adic} 的内容：任意 canonical finite-time-depth prime-register 机制都可忠实嵌入
\[
\mathrm{Aff}\!\bigl(T_2(E_{\min})\bigr),
\qquad
T_2(E_{\min})\cong \ZZ_2^2.
\]
第二条由定理 \ref{thm:xi-localized-phase-sampling-closed-form} 与定理 \ref{thm:conclusion-finite-prime-localization-pontryagin-rigidity} 给出：局部化相位采样只能唯一恢复有限素数集 \(S\)，等价地只能恢复
\[
0\to \prod_{p\in S}\ZZ_p\to \Sigma_S\to \TT\to 0
\]
中的 prime-ledger 核。第三条则是定理 \ref{thm:killo-localization-circle-dimension-prime-ledger-splitting} 的直接内容。

最后，由定理 \ref{thm:conclusion-finite-prime-solenoid-terminal-object}，\(\Sigma_S\) 的任意非平凡连通紧 Lie 商都同构于 \(\TT\)。因此二维 \(2\)-进实现空间并不意味着结构层存在第二个连通圆方向；它只是在 faithful implementation 中用于承载仿射动力学的 ambient coordinate。证毕。
\end{proof}

\begin{theorem}[有限素数支撑包含关系的商映射判别]\label{thm:conclusion-finite-primesupport-quotient-order-characterization}
对任意有限素数集 \(S,T\subset \mathbb P\)，下列命题等价：
\begin{enumerate}
  \item \(S\subseteq T\)；
  \item 存在连续满射
  \[
  q:\Sigma_T\twoheadrightarrow \Sigma_S
  \]
  ，其核全不连通，并且在标准圆商上满足
  \[
  \pi_S\circ q=\pi_T,
  \]
  其中 \(\pi_U:\Sigma_U\twoheadrightarrow \TT\) 是定理 \ref{thm:conclusion-finite-prime-solenoid-terminal-object} 中的标准投影；
  \item 存在离散侧单射
  \[
  \iota:G_S\hookrightarrow G_T.
  \]
\end{enumerate}
因此，有限素数支撑的包含关系，可以完全在 compact 完成侧用“是否存在带全不连通核的标准商映射”来刻画。
\end{theorem}
\begin{proof}
\(1\Rightarrow 2\) 正是定理 \ref{thm:conclusion-finite-primesupport-exact-quotient}，其中 \(q=q_{T\to S}\)。

\(2\Rightarrow 3\) 由 Pontryagin 对偶直接得到：连续满射 \(q\) 的对偶是离散侧单射
\[
\widehat q:G_S\hookrightarrow G_T.
\]

现证 \(3\Rightarrow 1\)。设 \(\iota:G_S\hookrightarrow G_T\) 为单射。若 \(p\in S\)，则 \(1\in G_S\) 是 \(p\)-可除的，因为
\[
1=p\cdot \frac{1}{p}
\qquad\text{且}\qquad
\frac{1}{p}\in G_S.
\]
更强地，对每个 \(n\ge 1\) 都有
\[
1=p^n\cdot \frac{1}{p^n},
\qquad
\frac{1}{p^n}\in G_S,
\]
故 \(\iota(1)\neq 0\) 在 \(G_T\) 中对一切 \(n\) 都是 \(p^n\)-可除的。若 \(p\notin T\)，则 \(G_T=\ZZ[T^{-1}]\) 中不存在非零的无限 \(p\)-可除元素：任取
\[
x=\frac{a}{\prod_{q\in T}q^{k_q}}\neq 0,
\]
若对每个 \(n\) 都存在 \(y_n\in G_T\) 使 \(x=p^n y_n\)，则分子 \(a\) 必被任意高次 \(p\) 整除，只能推出 \(a=0\)，与 \(x\neq 0\) 矛盾。因此 \(p\notin T\) 不可能。故 \(p\in T\)，从而 \(S\subseteq T\)。证毕。
\end{proof}

\begin{corollary}[有限素数支撑完成族的完全 Boolean 实现]\label{cor:conclusion-finite-primesupport-boolean-exact-realization}
赋值
\[
S\longmapsto \Sigma_S
\]
把有限素数支撑的 Boolean 格
\[
(\mathcal P_f(\mathbb P),\cup,\cap)
\]
完全实现到 compact abelian groups 范畴中，并满足
\[
\Sigma_{S\cup T}
\cong
\Sigma_S\times_{\Sigma_{S\cap T}}\Sigma_T,
\]
\[
S\subseteq T
\iff
\exists\,\Sigma_T\twoheadrightarrow \Sigma_S
\text{ with totally disconnected kernel}.
\]
因此，项目中的有限 prime-stage 已不再只是标签集合，而是一套可在 exact 范畴语义下完整演算的 Boolean support calculus。
\end{corollary}
\begin{proof}
第一条正是定理 \ref{thm:conclusion-finite-primesupport-bicartesian-boolean-calculus} 对 kernel 的 fiber-product 解释；第二条正是定理 \ref{thm:conclusion-finite-primesupport-quotient-order-characterization}。证毕。
\end{proof}

\begin{conjecture}[一轴 prime-register 的有限支撑逆极限骨架]\label{conj:conclusion-finite-primesupport-universal-inverse-limit-solenoid}
沿有限素数支撑格的包含关系，以定理 \ref{thm:conclusion-finite-primesupport-exact-quotient} 的过渡映射
\[
q_{T\to S}:\Sigma_T\twoheadrightarrow \Sigma_S
\qquad (S\subseteq T)
\]
组织出的 inverse system
\[
\bigl(\Sigma_S,q_{T\to S}\bigr)_{S\in \mathcal P_f(\mathbb P)}
\]
应当给出猜想 \ref{conj:conclusion-uniaxial-prime-register-universality} 中一轴 prime-register 普适对象的有限支撑骨架；更精确地，存在规范同构
\[
\Sigma_{\mathrm{univ}}
\cong
\varprojlim_{S\in \mathcal P_f(\mathbb P)}\Sigma_S.
\]
若此猜想成立，则：
\begin{enumerate}
  \item 定理 \ref{thm:conclusion-finite-primesupport-exact-quotient}--\ref{thm:conclusion-finite-primesupport-quotient-order-characterization} 所给出的 finite-support Boolean exact calculus 将汇聚为一个真正的 universal one-axis phase register；
  \item 全部有限 prime-stage completions 都将成为该一轴对象的 cofinal 商；
  \item “有限状态不可逆性可由单轴 prime-register 逆极限统一吸收”将由有限支撑近似提升为严格的 inverse-limit theorem。
\end{enumerate}
\end{conjecture}

上述链条把有限素数局部化、局部化相位采样与 canonical prime-register 的 faithful implementation 压成了同一骨架：有限素数支撑在 compact--profinite 侧并非一组松散标签，而是携带精确商、精确交与精确粘接的 Boolean exact calculus；而一轴 prime-register 的普适性，则应当正是这套 finite-support exact calculus 的 cofinal inverse-limit completion。
